\documentclass[11pt]{article}
\usepackage{fontspec}
\usepackage[french,english]{babel}
\usepackage{amsmath}
\usepackage{amssymb}
\usepackage{amsthm}
\usepackage{mathrsfs}
\usepackage{tikz-cd}
\usepackage{enumitem}

\frenchsetup{AutoSpacePunctuation=false}

\newcommand{\printedpage}[1]{\marginpar{\footnotesize p.~#1}}
\newcommand{\msp}[1]{\marginpar{\footnotesize [#1]}}
\newcommand{\Ch}{\underline{\mathrm{Ch}}}

\begin{document}

\printedpage{9}
% head: A. Grothendieck : Dérivateurs, Ch. XII, Édition des Universités de Montpellier II et Paris VII

Thus, the lemma on page 7 is proved, using axiom W(4) which ensures that an arrow which is both $W$-smooth and $W$-proper satisfies the very strong condition a) of $W$-fibrations.

\textbf{NB} In fact, it would suffice, instead of axiom W(4) in its strong form, to assume that an $f : X \longrightarrow Y$ which is both ``homotopically smooth'' and ``homotopically proper'', i.e. the fibers of the $X/x \longrightarrow Y/y$ and of the $x\backslash X \longrightarrow y\backslash Y$ are contractible (hence $W$-aspheric for every $W$ satisfying W(1,2,3)), is a strong $W$-fibration. Let us state all the same what we have proved.

\textbf{Proposition 4.} \textit{If $W$ satisfies} W(1,2,3) (\footnote[17]{For W(2), it suffices in the form W(2 bis).}) (we always assume this in what follows) \textit{plus} W(4) (in a somewhat weakened sufficient form, cf. preceding remark), \textit{we have this: Let $f : X \longrightarrow Y$ be an arrow in} Cat$/S$\textit{, with $X$, $Y$ $W$-fibered over $S$. If $f \in W$, then $f_s : X_s \longrightarrow Y_s$ is in $W$ for every $s \in S$.}

I wonder whether I know enough to be able to prove
\[ W \subseteq W_\infty . \]
We must all the same eliminate joke cases, such as $W = \mathrm{Fl}(\mathrm{Cat})$, hence put in the axiom

\begin{itemize}[leftmargin=4em, labelwidth=3em, labelsep=0.6em]
\item[W(5)] (Axiom of connectedness (\footnote[18]{See on this subject the corollary on page 22.}).) If $f : X \longrightarrow Y$ is in $W$, $\pi_0(f) : \pi_0(X) \longrightarrow \pi_0(Y)$ \texttt{[is]} bijective.
\end{itemize}

\msp{page 11}I then factor $f$ as (\footnote[19]{Once again, in these arguments it pays to work with $W_\omega$.})
\[ X \xrightarrow{i} \bar{X} \xrightarrow{\bar{f}} Y, \]
with $\bar{f}$ a $W_\infty$-fibration and a $W$-fibration (\footnote[20]{in fact, $\bar{f}$ is both $W_\infty$-perfect (i.e. $W_\infty$-smooth and proper) and $W$-perfect (in fact, $W_\omega$-perfect).}), $i \in W \cap W_\infty$ - this works after all, % [sic] finallement -> finalement, spelling
by the construction in VII, p. 170-171 (of \textsc{Cartan-Serre}), starting from the properties of % [sic] propiétés -> propriétés, missing letter
$Y \longrightarrow \Ch(Y) \longrightarrow Y \times Y$ and nothing more. As $\bar{f} i = f \in W$ and $i \in W$, we deduce $\bar{f} \in W$. As $i \in W_\infty$, we conclude $f \in W_\infty \iff \bar{f} \in W_\infty$. Thus to prove that $f$ is in $W_\infty$, we are reduced to the case where $f$ is a $W$-fibration and a $W_\infty$-fibration. But then proposition 4 shows us that the fibers of $f$ are $W$-aspheric (\footnote[21]{we have used here only this property of the fibers, as a consequence of W(4). (We use W(4) in its weakest form W(4b$''$), cf. pages 13, 14).}), and it suffices to see that they are $W_\infty$-aspheric (proposition 3 (\footnote[22]{we use proposition 3 only in the case $Y = S$, for which W(2 bis) suffices.}) applied to $W_\infty$). Thus we are reduced to the case where $Y = e$, i.e. to proving that if $X$ in Cat is $W$-aspheric, it is $W_\infty$-aspheric.
% [sic] (footnote 21) consequence -> conséquence, missing accent

Let $x_0 \in \mathrm{Ob}\, X$, let us consider as before
\[ E = \Ch(X, x_0, -) \xrightarrow{p} X, \]
% foot: Version 7 décembre 2002
\printedpage{10}
% head: A. Grothendieck : Dérivateurs, Ch. XII, Édition des Universités de Montpellier II et Paris VII
we know that $E$ is $W \cap W_\infty$-aspheric, and that it is a fibration for $W$ and for $W_\infty$. As $X$ is $W$-aspheric (like $E$), $p \in W$, hence again by proposition 4 (\footnote[23]{in fact, condition W(4b$''$) on $W$.}), the fibers of $p$ are $W$-aspheric, in particular (by W(5)) 0-connected. We have therefore proved the

\msp{page 12}\textbf{Lemma.} (\footnote[24]{Uses W(1, 2 bis, 3, 4b$''$) on $W$.}) \textit{Let $(X, x_0)$ be in} Cat$_\bullet$\textit{. If $X$ is $W$-aspheric, then $\Omega(X, x_0) \overset{\text{def}}{=} \Ch(X; x_0, x_0)$ is too} (assuming axioms W(1,2,3,4) on $W$)\textit{. Thus (by iteration) the $\Omega^n(X, x_0)$ are $W$-aspheric} ($\forall\, n \in \mathbf{N}$).

\textbf{Corollary.} \textit{If moreover $W$ satisfies the axiom} W(5) \textit{of connectedness, then $X$ is $W_\infty$-aspheric.}

Indeed, ``we know'' that this means that all the $\Omega^n(X, x_0)$ are 0-connected. Now the connectedness axiom W(5) ensures that $W$-aspheric implies 0-connected.

But the ``we know'' is a tiny bit flimsy (\footnote[25]{This is proved quite simply in the corollary to ``theorem 1'', page 190.}). To justify it, it would be necessary, either % [sic] (footnote 25) theorème -> théorème, missing accent
that I make the link with classical 1/2-simplicial theory, and prove that the $\Omega^n$ which I construct are ``the right ones'', i.e. that their nerve does indeed give back (up to isomorphism in Hot) the classical $\Omega^n$ (??). In particular, that the $\pi_0(\Omega^n(X, x_0))$ is indeed ``everybody's'' $\pi_n(X, x_0)$. Or else that I prove directly that the 0-connectedness of the $\Omega^n(X, x_0)$ is equivalent to the $W_\infty$-asphericity of $X$, defined cohomologically in the manner of \textsc{Artin-Mazur}. So this is a piece of work that I shall absolutely have to do. But it is now no more than a simple piece of routine work, and I can state right now the

\textbf{Theorem 1.} (\footnote[26]{It even suffices to have W(4) in a weakened form, with ``fibrant'' functors in an ``absolute'' sense. Cf. the NB p. 10 - and above all, cf. pages 13, 14.},\footnote[27]{NB Instead of W(5), it suffices to assume $W \neq \mathrm{Fl}(\mathrm{Cat})$, cf. th. 2, p. 20.},\footnote[28]{The useful hypotheses on $W$ are in the end W(1), W(2 bis), W(3), W(4b$''$), and $W \neq \mathrm{Fl}(\mathrm{Cat})$.}) \textit{Let $W \subseteq$} Fl(Cat) \textit{satisfy the conditions} W(1,2,3,4,5) (axiom of saturation, of localization, of the final object, of fibration, of connectedness)\textit{. Then}
\[ W \subseteq W_\infty . \]

\textbf{Corollary.} \textit{We can characterize $W_\infty$ as the largest subset of} Fl(Cat) \textit{satisfying the}
\msp{page 13}\textit{conditions} W(1) \textit{to} W(5) \textit{above.}

Of course, one would like a more interesting axiom 6, which joined to the preceding ones ensures that $W = W_\infty$ (\footnote[29]{Cf. axiom W(6) on the next page.}). Or do these properties already ensure that $W = W_\infty$? To see this, perhaps I would have to know more about the structure of $W_\infty$-aspheric objects? If this were false, I am thinking of an axiom of ``homotopy descent'', which would imply that a $W_\infty$-aspheric object is also $W$-aspheric.

% foot: Version 7 décembre 2002
\printedpage{11}
% head: A. Grothendieck : Dérivateurs, Ch. XII, Édition des Universités de Montpellier II et Paris VII

I note right now that if some $W_i$ ($i \in I$) satisfy axioms W(1) to W(5), the same is true of $\bigcap_i W_i$ (\footnote[30]{This is true for any set of axioms in the nature of stability axioms.}). Thus there is a smallest set of arrows $W_\omega$ satisfying these conditions. If it is not equal to $W_\infty$, i.e. if axioms W(1) to W(5) do not already characterize $W_\infty$, it would be a really very interesting localizer.

I realize that I have used only a rather weak form of axiom W(4) - essentially, the validity of proposition 4. I am going to summarize different variants of the fibration axiom.

\textbf{W(4)}

\begin{itemize}[leftmargin=4em, labelwidth=2.5em, labelsep=0.6em, align=right]
\item[a$_1$)] If $f : X \longrightarrow Y$ in Cat is a left $W$-fibration, it is an ultrastrong $W$-fibration.

\item[a$_2$)] If $f : X \longrightarrow Y$ is a right $W$-fibration, it is an ultrastrong $W$-fibration.

\item[a$'$)] If $f : X \longrightarrow Y$ is $W$-perfect, i.e. both $W$-proper and $W$-smooth (\footnote[31]{which implies that it is a left $W$-fibration \textit{and} a right one.}), then $f$ is an ultrastrong $W$-fibration.

\item[a$''$)] If $f : X \longrightarrow Y$ is both homotopically proper and homotopically smooth (cf. NB page 10), $f$ is an ultrastrong $W$-fibration. (It is better to say: If $f$ is $W_\omega$-perfect, it is an ultrastrong $W$-fibration.)

\item[b)] Let $f : X \longrightarrow Y$ be an arrow in Cat$/S$, with $X$ and $Y$ $W$-fibered over $S$. If $f \in W$, then $\forall\, s \in S$, $f_s : X_s \longrightarrow Y_s$ is \texttt{[in]} $W$.

\item[b$'$)] Let $f : X \longrightarrow Y$ be an arrow of Cat$/S$ which is $W$-perfect and $W_\infty$-perfect, then $f \in W$ implies that the fibers of $f$ are $W$-aspheric. (It is better to say: If $f$ is $W_\omega$-perfect, it is an ultrastrong $W$-fibration.) \texttt{[Rather: Let $f : X \longrightarrow Y$ be an arrow in \textrm{Cat}$/S$ such that $X \longrightarrow S$ and $Y \longrightarrow S$ are $W_\omega$-perfect. If $f \in W$, then $\forall\, s \in S$, $f_s : X_s \longrightarrow Y_s$ is in $W$.]}

\item[b$''$)] \msp{page 14}Let $W_\omega \subseteq W \cap W_\infty$ be the smallest subset of Fl(Cat) satisfying conditions W(1,2,3) (\footnote[32]{It is better to take W(2 bis) instead of W(2), to find $W_\omega$ as small as possible!}) (axioms of saturation, of localization, of the final object). Let $f : X \longrightarrow Y$ be an arrow in Cat which is $W_\omega$-perfect. If $f \in W$, then the fibers of $f$ are $W_\omega$-aspheric \texttt{[rather $W$-aspheric]}.
\end{itemize}

Among these properties of $W$, we have the tautological implications (except a$''$) $\Longrightarrow$ b), which is essentially proposition 4)

% foot: Version 7 décembre 2002
\printedpage{12}
% head: A. Grothendieck : Dérivateurs, Ch. XII, Édition des Universités de Montpellier II et Paris VII

\[
\begin{tikzcd}[row sep=huge, column sep=0pt]
\mathrm{a} \arrow[d, Rightarrow] & {}= (\mathrm{a}_1 \Longleftrightarrow \mathrm{a}_2) & \hspace{2.5cm} \\
\mathrm{a}' \arrow[d, Rightarrow] & & \smash{\left.\rule[-1.8cm]{0pt}{3.6cm}\right\}\ \text{\begin{tabular}{l}In fact,\\ all these\\ conditions are\\ equivalent\end{tabular}}} \\
\mathrm{a}'' \arrow[d, Rightarrow] & & \\
\mathrm{b} \arrow[d, Rightarrow] & & \\
\mathrm{b}' \arrow[d, Rightarrow] & & \\
\mathrm{b}'' & &
\end{tikzcd}
\]

(\footnote[33]{I return to the various variants of W(4) below, p. 53-56.}).

Well, in the lemma on page 12, it sufficed to assume that $W$ satisfies the weakest version b$''$) of the axiom of fibrations (seven versions!). Indeed, the fibrations $\bar{X} \longrightarrow Y$ and $E \longrightarrow X$ used in the proof of the lemma are $W_\omega$-perfect, and it sufficed to know that if they are in $W$, then the fibers are $W$-aspheric.

But I do not think that even the strongest version W(4a) of the axiom of fibrations, joined to the three preceding axioms and to W(5) (connectedness), implies $W = W_\infty$. It would at least be necessary to add to it the axiom

\begin{itemize}[leftmargin=4em, labelwidth=3em, labelsep=0.6em]
\item[\textbf{W(6)}] (Axiom of filtering limits.) Let $X = \varinjlim\limits_{i} X_i$ be a filtering direct limit in Cat. If the $X_i$ are $W$-aspheric, $X$ too.
\end{itemize}

This should be equivalent to the property, apparently stronger, in terms of arrows of $W$:

\textbf{Proposition 5.} \textit{Suppose we have} W(6) (in addition to W(1) to W(4) (\footnote[34]{W(4) in the weakest form W(4b$''$) .})),
\msp{page 15}\textit{and let}
\[ (f_i)_{i\in I} : (X_i)_{i\in I} \longrightarrow (Y_i)_{i\in I} \]
% foot: Version 7 décembre 2002
\printedpage{13}
% head: A. Grothendieck : Dérivateurs, Ch. XII, Édition des Universités de Montpellier II et Paris VII
\textit{a homomorphism of direct systems in} Cat\textit{, indexed by a filtering ordered set, let $X = \varinjlim X_i$, $Y = \varinjlim Y_i$, $f = \varinjlim f_i$, $f : X \longrightarrow Y$. If $f_i$ is in $W$ for every $i$, then $f \in W$.}

We factor each $f_i$ as
\[ X_i \xrightarrow{\alpha_i} \bar{X}_i \xrightarrow{\bar{f}_i} Y_i, \qquad \alpha_i \in W_\omega,\ \bar{f}_i \in \mathrm{Parf}\, W_\omega \]
(\footnote[35]{$W_\omega$ introduced in condition W(4b$''$), page 14.}) following the patented system. The construction being functorial, the $\bar{X}_i$ too form a direct system. Moreover, as the functor $S \longrightarrow \Ch(S)$ which gives rise to the canonical factorization of morphisms, commutes with filtering $\varinjlim$ (immediate), it follows that, if we put
\[ \bar{X} = \varinjlim_{i} \bar{X}_i \]
and look at the arrows
\[ X \xrightarrow{\alpha} \bar{X} \xrightarrow{\bar{f}} Y, \qquad \alpha = \varinjlim \alpha_i,\ \bar{f} = \varinjlim \bar{f}_i, \]
with $\bar{f}\alpha = \bar{f}$, this factorization of $f$ is indeed its canonical factorization. Thus $\alpha \in W$, % [sic] $\bar{f}\alpha = \bar{f}$ -> $\bar{f}\alpha = f$, the composite should be f
$\bar{f} \in \mathrm{Parf}\, W_\omega$. This having been seen, by virtue of prop. 3, $f \in W$, i.e. $\bar{f} \in W$, is true if the fibers of $\bar{f}$ are $W$-aspheric. But the fibers of $\bar{f}$ are interpreted as filtering direct limits of fibers of the $\bar{f}_i$, which are $W$-aspheric by condition W(4b$''$) on fibrations. Thus axiom W(6) gives what is needed.

\textbf{Conjecture} (bold? foolish?). \textit{If $W \subseteq$} Fl(Cat) \textit{satisfies the axioms} W(1) \textit{to} W(6) \textit{(with} W(4) \textit{in its weakest form} W(4b$''$) (\footnote[36]{and perhaps W(2) in the weaker form W(2 bis).})\textit{), then $W = W_\infty$.}

\msp{page 16}But it is quite possible that one must strengthen W(1) (saturation) a little, by requiring that every retract of an $f \in W$ be in $W$.

\textbf{Remark.} Let us consider, for $n \in \mathbf{N}$, the set $W_n \subseteq$ Fl(Cat) of ``$n$-homotopy equivalences'' in Cat:
\[ f \in W_n \iff \pi_0(f) \text{ bijective, and } \pi_i(f, x) : \pi_i(X, x) \longrightarrow \pi_i(Y, f(x)) \text{ bijective for } 0 < i \leq n. \]

I thought that this satisfied all that one can ask of a localizer, without having taken the trouble to write it down carefully. But this cannot satisfy W(1) through W(4b$''$), since we % [sic] Ma -> Mais, word truncated
do not have $W_n \subseteq W_\infty$!

W(1) is trivial from the definition (and we have without tears even the \textit{strong saturation} of $W_i$, % [sic] $W_i$ -> $W_n$ (also in the following lines), the remark concerns W_n
i.e. if $f$ in Cat is \texttt{[an]} isomorphism in $\mathrm{Hot}_{W_i} = (\mathrm{Cat})W_i^{-1}$, then $f \in W_i$). I think I remember having checked W(2), by the same method as for $W_\infty$, using a cohomological criterion in the manner of \textsc{Artin-Mazur}, to express $f \in W_i$, and using the Leray spectral sequences for $X \longrightarrow S$ and $Y \longrightarrow S$. Condition W(3) is trivially satisfied since $W_n \supseteq W_\infty$ and it is for $W_\infty$. So it gets botched where I
% foot: Version 7 décembre 2002

\end{document}
